\section{DALTON Dataset and Problem Formulation}
\label{sec:dataset}

\subsection{Dataset Overview and Characteristics}

The DALTON indoor air quality dataset comprises 89.1 million samples collected from 30 residential sites across India between October 2022 and March 2023, spanning approximately 13,646 cumulative hours of continuous monitoring \cite{karmakar2024indoor}. Each DALTON node measures seven pollutant channels at 1 Hz sampling frequency: particulate matter concentrations (PM$_{2.5}$, PM$_{10}$ in $\mu$g/m$^3$), carbon dioxide (CO$_2$ in ppm), nitrogen dioxide (NO$_2$ in ppm), total volatile organic compounds (VOC in ppb), ethanol (in ppm), and carbon monoxide (CO in ppm). Sites encompass diverse geographic and socioeconomic contexts: 12 rural villages (mud-brick homes with biomass cooking), 10 suburban towns (concrete houses with LPG stoves), and 8 urban slums (multi-story apartments with mixed ventilation) \cite{karmakar2024exploring}. This diversity enables evaluation of model generalizability across building types, fuel sources, and occupant behaviors.

Figure~\ref{fig:dalton_timeseries} illustrates representative 24-hour time series from three sites, highlighting characteristic patterns: (i) diurnal cycles with morning (6-9 AM) and evening (6-10 PM) pollution peaks coinciding with cooking activities, (ii) quasi-static baseline levels during nighttime and daytime occupied periods, and (iii) transient spikes during cooking events (PM$_{2.5}$ increases from $\sim$30 to $>$400 $\mu$g/m$^3$ within 10 minutes). Autocorrelation analysis reveals that pollutant concentrations exhibit Pearson correlation $\rho > 0.95$ for lags up to 5 minutes, decaying to $\rho \approx 0.6$ at 30-minute lags, confirming the slow temporal dynamics that motivate residual encoding (Section~\ref{sec:methodology}).

\subsection{Site Selection and Data Preprocessing}

For this study, we select 10 representative sites stratified by geography (4 rural, 3 suburban, 3 urban) and seasonal coverage (sites with $>$400 hours spanning both winter and summer). This subset contains 29.7 million samples, sufficient for training deep models while maintaining computational tractability. Data preprocessing follows these steps:

\begin{enumerate}
\item \textbf{Outlier removal}: Eliminate intervals where sensor faults are flagged in DALTON metadata (e.g., calibration periods, power failures), accounting for 2.3\% of samples.
\item \textbf{Gap interpolation}: For short gaps ($<$60 seconds), apply linear interpolation; discard longer gaps to preserve temporal continuity.
\item \textbf{Normalization}: Compute per-site, per-channel mean $\mu_c^{(s)}$ and standard deviation $\sigma_c^{(s)}$ on training data, then normalize as $z_t^{(c)} = (x_t^{(c)} - \mu_c^{(s)}) / \sigma_c^{(s)}$. This accounts for site-specific baseline concentrations (e.g., rural sites have lower PM$_{10}$ baselines than urban sites).
\item \textbf{Temporal windowing}: Construct overlapping sliding windows of length $L = 256$ seconds (256 samples at 1 Hz) with stride 64 seconds, yielding $\sim$1.8 million training windows.
\end{enumerate}

Normalization parameters $\{\mu_c^{(s)}, \sigma_c^{(s)}\}$ are stored at the edge gateway for denormalization of TCN predictions, incurring negligible storage overhead (12 floats per site).

\subsection{Problem Formulation: Forecasting Task}

\subsubsection{Input and Target Definition}

We define a supervised regression task aligned with predictive IAQ control applications \cite{quang2025ai}. At time $t$, the model receives a multivariate time series of quantized pollutant measurements:
\begin{equation}
\mathbf{X}_{t-L:t} = \left[ \tilde{\mathbf{x}}_{t-L+1}, \tilde{\mathbf{x}}_{t-L+2}, \ldots, \tilde{\mathbf{x}}_t \right] \in \mathbb{R}^{L \times C},
\label{eq:input_sequence}
\end{equation}
where $\tilde{\mathbf{x}}_t = [\tilde{x}_t^{(1)}, \ldots, \tilde{x}_t^{(C)}]^\top$ is the quantized observation vector at time $t$, $L = 256$ is the window length, and $C = 6$ is the number of pollutant channels (we exclude ethanol, which exhibits low variability and high correlation with VOC \cite{karmakar2024exploring}).

The target is the average PM$_{2.5}$ concentration over a 10-minute horizon $[t+1, t+600]$:
\begin{equation}
y_t = \frac{1}{600} \sum_{k=1}^{600} x_{t+k}^{(\text{PM}_{2.5})}.
\label{eq:target}
\end{equation}

This formulation supports proactive interventions (e.g., activating ventilation, alerting occupants) with sufficient lead time, consistent with prior IAQ forecasting studies \cite{nguyen2025aiot, yildiz2025development}.

\subsubsection{Objective Function}

The TCN model $f_\theta: \mathbb{R}^{L \times C} \rightarrow \mathbb{R}$ parameterized by weights $\theta$ minimizes the mean squared error (MSE) over training windows:
\begin{equation}
\mathcal{L}(\theta) = \frac{1}{N} \sum_{i=1}^{N} \left( y_i - f_\theta(\mathbf{X}_i) \right)^2,
\label{eq:objective}
\end{equation}
where $N$ is the number of training windows. During inference, predictions are evaluated using root mean squared error (RMSE) and mean absolute error (MAE):
\begin{equation}
\text{RMSE} = \sqrt{\frac{1}{M} \sum_{j=1}^{M} (y_j - \hat{y}_j)^2}, \quad
\text{MAE} = \frac{1}{M} \sum_{j=1}^{M} |y_j - \hat{y}_j|,
\label{eq:metrics}
\end{equation}
where $M$ is the number of test windows and $\hat{y}_j = f_\theta(\mathbf{X}_j)$ is the predicted target.

\subsection{Train-Validation-Test Splits}

To respect temporal ordering and prevent data leakage, we employ chronological splits at each site independently. For a site with total duration $T$ seconds:
\begin{itemize}
\item \textbf{Training set}: First 60\% of chronological data ($t \in [0, 0.6T]$).
\item \textbf{Validation set}: Next 20\% ($t \in [0.6T, 0.8T]$), used for hyperparameter tuning and early stopping.
\item \textbf{Test set}: Final 20\% ($t \in [0.8T, T]$), held out for final performance evaluation.
\end{itemize}

This split ensures that test windows represent future time periods unseen during training, mimicking real-world deployment scenarios. For sites spanning multiple months, test sets include both winter and summer conditions, enabling assessment of seasonal generalization.

\subsection{Temporal Dynamics and Autocorrelation Analysis}

To justify design choices in Section~\ref{sec:methodology} (e.g., residual encoding, subsampling interval), we analyze DALTON's temporal characteristics. Figure~\ref{fig:autocorrelation} plots autocorrelation functions (ACF) for PM$_{2.5}$, CO$_2$, and VOC averaged across 10 sites:
\begin{equation}
\text{ACF}(\tau) = \frac{\mathbb{E}[(x_t - \mu)(x_{t+\tau} - \mu)]}{\sigma^2},
\label{eq:acf}
\end{equation}
where $\tau$ is the lag in seconds, $\mu$ and $\sigma^2$ are the mean and variance of the time series.

Key observations:
\begin{itemize}
\item PM$_{2.5}$ and CO$_2$ maintain $\text{ACF} > 0.95$ for lags $\tau < 300$ s (5 minutes), indicating strong short-term predictability.
\item VOC and NO$_2$ exhibit faster decorrelation ($\text{ACF} \approx 0.7$ at $\tau = 300$ s), reflecting transient emission events.
\item At $\tau = 240$ s (4 minutes), all pollutants retain $\text{ACF} > 0.85$, justifying $T_S = 4$ s subsampling without significant loss of temporal structure.
\end{itemize}

Power spectral density (PSD) analysis (not shown) reveals dominant frequencies at $f < 0.01$ Hz (periods $>$ 100 s), confirming that high-frequency components ($f > 0.1$ Hz, periods $<$ 10 s) primarily represent sensor noise rather than true pollutant dynamics. This supports aggressive temporal subsampling and quantization without degrading signal-to-noise ratio.

\subsection{Summary}

The DALTON dataset provides a realistic, large-scale testbed for validating energy-efficient air quality monitoring schemes. Its temporal characteristics—slow dynamics, high autocorrelation, and bounded ranges—align with the design principles of adaptive quantization and residual encoding proposed in Section~\ref{sec:methodology}. The next section describes experimental protocols and baseline comparisons.